\documentclass[12pt,leqno]{article}

%----------------------------------------------------------------
% 1. DOCUMENT SETUP & TYPOGRAPHY
%----------------------------------------------------------------
\usepackage{math}
\usepackage[TS1, T1]{fontenc}               % Font encoding for output
\usepackage[adobe-utopia]{mathdesign}       % Adobe Utopia font
\usepackage{microtype}                      % Improves typography (line breaking, spacing)
\usepackage{setspace}                       % For line spacing control (e.g., \doublespacing)
\usepackage{soul}                           % For strikethrough and underlining (\ul)
\usepackage{dirtytalk}                      % For inline quotations
\usepackage{titling}                        % Control over \maketitle
\usepackage{epigraph}                       % For epigraphs
\usepackage{booktabs}                       % Required for importing nice tables
\usepackage{tcolorbox}                      % Colour box
\usepackage[ruled,vlined]{algorithm2e}
\usepackage{float}

\tcbuselibrary{breakable}                   % Required for page breaking

%----------------------------------------------------------------
% 2. PAGE LAYOUT & LINKS
%----------------------------------------------------------------
\usepackage{geometry}                       % For page margins and layout
\usepackage{pdflscape}                      % For landscape pages
\usepackage[dvipsnames]{xcolor}             % Color support
\usepackage[colorlinks=true]{hyperref}      % PDF hyperlinks (should be loaded late)

%----------------------------------------------------------------
% 3. MATHS & THEOREMS
%----------------------------------------------------------------
\usepackage{amsmath}                        % Core maths tools

% Theorem-like environments
\newtheorem{theorem}{Theorem}[section]
\newtheorem{assumption}{Assumption}
\newtheorem{corollary}[theorem]{Corollary}
\newtheorem{proposition}{Proposition}[section]
\newtheorem{lemma}[theorem]{Lemma}
\newtheorem{conjecture}{Conjecture}[section]
\newenvironment{proof}[1][Proof]{\noindent\textbf{#1.} }{\ \rule{0.5em}{0.5em}}

%----------------------------------------------------------------
% 4. TABLES, FIGURES & GRAPHICS
%----------------------------------------------------------------
\usepackage{graphicx}                       % For including images
\usepackage{caption, subcaption}            % Enhanced captions for figures and tables
\usepackage{booktabs}                       % High-quality table rules (\toprule, \midrule, \bottomrule)
\usepackage{dcolumn}                        % Align table columns on a decimal point
\usepackage{longtable, threeparttable, threeparttablex} % For long tables and tables with notes
\usepackage{multirow, adjustbox}            % For multi-row cells and adjusting box content
\usepackage{rotating}                       % For rotating elements (e.g., sidewaystable)
\usepackage{tikz}                           % For drawing graphics
\usetikzlibrary{positioning}

%----------------------------------------------------------------
% 5. LISTS
%----------------------------------------------------------------
\usepackage{enumitem}                       % Advanced control over lists

%----------------------------------------------------------------
% 6. BIBLIOGRAPHY (Author-Year Style)
%----------------------------------------------------------------
\usepackage[
    backend=biber,      % Use the Biber engine
    style=authoryear,   % Set main style to author-year
    citestyle=authoryear-comp, % Compress citations (e.g., Author 2020, 2021)
    natbib=true,        % Allow use of natbib commands like \citet and \citep
    maxcitenames=2,     % Max authors to show in citation before "et al."
    maxbibnames=99,     % Max authors to show in the bibliography
    uniquelist=false    % Prevents adding initials to distinguish authors
]{biblatex}
\addbibresource{bibliography.bib}           % Specify your .bib file

%----------------------------------------------------------------
% 7. GLOBAL CONFIGURATIONS & SETTINGS
%----------------------------------------------------------------
% Page layout & colours
\geometry{left=1.25in, right=1.25in, top=1.25in, bottom=1.25in}
\definecolor{maroon}{rgb}{0.5, 0.0, 0.0}
\hypersetup{
    urlcolor   = maroon,
    linkcolor  = maroon,
    citecolor  = maroon 
}

% Typesetting adjustments
\setlength{\parindent}{0pt}
\interfootnotelinepenalty=10000        
\setlist[itemize]{noitemsep}

\begin{document}
%----------------------------------------------------------------
%% Title
%----------------------------------------------------------------
\title{\Large \bfseries \vspace{-5mm} PHD22 Macro labour: Problem set \vspace{-5mm}}
\author{Rob Włodarski}
\date{April 2026}
\maketitle

%------------------------------------------------------------------------
\section*{Question 1: The model}
I follow the exposition from \citet{hagedorn_cyclical_2008} wherever possible.
The Bellman equations that govern the model are:
\begin{subequations}
    \begin{align}
        & J\bp{p} = p - w\bp{p} + \beta \bp{1-s}  \E_p \bs{J\bp{p'}}, \\
        & V\bp{p} = -c + \beta q \bs{\theta\bp{p}}  \E_p \bs{J\bp{p'}}, \\ 
        & U\bp{p} = b + \beta \ba{f \bs{\theta\bp{p}} \E_p \bs{W\bp{p'}}  + \bc{1 - f \bs{\theta\bp{p}}}  \E_p \bs{U\bp{p'}}}, \\
        & W\bp{p} = w\bp{p} + \beta \bc{\bp{1-s} \E_p \bs{W\bp{p'}} + s \E_p \bs{U \bp{p'}} }.
    \end{align}
\end{subequations} 
I use the matching function from \citet{den_haan_job_2000} :
\begin{equation}
    m\bp{u,v}= \frac{uv}{\bp{u^l + v^l}^{\frac{1}{l}}},
\end{equation} 
which gives the following contact rates:
\begin{subequations}
    \begin{align}
        & q\bs{\theta\bp{p}} = \bs{1+\theta\bp{p}^{l}}^{-\frac{1}{l}}, \label{jfr}\\
        & f\bs{\theta\bp{p}} = \bs{1+\theta\bp{p}^{-l}}^{-\frac{1}{l}}.
    \end{align}
\end{subequations}

\begin{tcolorbox}[sharp corners,breakable, colback=gray!10, colframe=gray!90, title= Free entry condition]
    The \textbf{free entry condition} stipulates $V\bp{p} = 0$
    which translates to:
    \begin{subequations}
        \begin{align}
            & c = \beta q\bs{\theta\bp{p}}  \E_p \bs{J\bp{p'}}, \\
            & J\bp{p} = p - w\bp{p} + \frac{\bp{1-s}c}{ q\bs{\theta\bp{p}}}.
        \end{align}
    \end{subequations}
\end{tcolorbox}
The free entry condition can be used to derive the job posting condition. 
In the steady state, where we have only one price, $p=p^{\star}$, the conditions collapse to:
\begin{subequations}
        \begin{align}
            & c = \beta q\bp{\theta^{\star}}  J^{\star} \implies  J^{\star} = \frac{c}{\beta q\bp{\theta^{\star}} }\\
            & J^{\star} = p^{\star} - w^{\star} + \beta \bp{1-s} J^{\star}  \implies J^{\star} = \frac{p^{\star} - w^{\star}}{1- \beta \bp{1-s}}.
        \end{align}
\end{subequations}
\begin{tcolorbox}[sharp corners,breakable, colback=gray!10, colframe=gray!90, title= Steady state job creation condition]
    Setting $r \equiv \frac{1}{\beta}-1$, the \textbf{job creation condition} becomes:
    \begin{equation}
        w^\star=p^\star - \frac{r+s}{q\bp{\theta^\star}c}. \label{jobcreation}
    \end{equation}
\end{tcolorbox}

\begin{tcolorbox}[sharp corners,breakable, colback=gray!10, colframe=gray!90, title= Nash bargaining and surplus sharing]
    A worker and a firm split the surplus, with workers taking $\gamma$:
    \begin{subequations}
        \begin{align}
            & W\bp{p} = \gamma \bc{J\bp{p}+W\bp{p} - U\bp{p}}, \\
            & J\bp{p} = \bp{1-\gamma} \bc{J\bp{p}+W\bp{p} - U\bp{p}}.
        \end{align}
    \end{subequations}
\end{tcolorbox}
Using the surplus sharing condition, I solve for the closed-form solution of $w\bp{p}$:
\begin{subequations}
    \begin{align}
         \bs{W}: \quad &W\bp{p} - \beta \E_p \bs{W\bp{p'}} = w\bp{p} - \beta s \E_p \bs{W\bp{p'} - U\bp{p'}},\\
         \bs{U}: \quad & U\bp{p} - \beta \E_p \bs{U\bp{p'}} = b + \beta  f\bs{\theta\bp{p}} \E_p \bs{W\bp{p'} - U\bp{p'}},
    \end{align}
    which can be combined:
    \begin{align}
        & W\bp{p} - U\bp{p} = w\bp{p} - b + \beta \bc{1-s -f\bs{\theta\bp{p}}} \E_p \bs{W\bp{p'} - U\bp{p'}}.
    \end{align}
    Now, by Nash bargaining and free entry, we must have that:
    \begin{align}
        & W\bp{p} - U\bp{p} = \frac{\gamma}{1-\gamma}J\bp{p} = \frac{\gamma}{1-\gamma} \bs{p - w\bp{p} + \frac{\bp{1-s}c}{ q\bs{\theta\bp{p}}} },\\
        & \E_p \bs{W\bp{p'} - U\bp{p'}} = \frac{\gamma}{1-\gamma}\E_p  \bs{J\bp{p}}=\frac{\gamma}{1-\gamma}\frac{c}{\beta q\bs{\theta\bp{p}}}.
    \end{align}
    The last three conditions can be combined and rearranged to solve for the wage.
\end{subequations}
\begin{tcolorbox}[sharp corners,breakable, colback=gray!10, colframe=gray!90, title= Equilibrium wage]
    By combining the free entry condition with Nash bargaining, I obtain the closed form wage:
    \begin{equation}
        w\bp{p} = \gamma \bs{p+\theta\bp{p}c} + \bp{1-\gamma}b. \label{eq_wage}
    \end{equation}
\end{tcolorbox}
\begin{tcolorbox}[sharp corners,breakable, colback=gray!10, colframe=gray!90, title= Beveridge curve]
    From the law of motion of employment: 
    \begin{subequations}
        \begin{align}
            & n' = \bp{1-s}n + m\bp{v,u}
        \end{align}
        I obtain the steady state Beveridge curve:
        \begin{align}
            u^\star = \frac{s}{f+s}. \label{beveridge}
        \end{align}
    \end{subequations}
\end{tcolorbox}
%------------------------------------------------------------------------
\section*{Question 2: The steady state algorithm}

\begin{algorithm}[H]
\caption{Calibration of vacancy posting cost to target unemployment} 
\KwIn{Parameters $(l,\gamma,b,\beta,s,p^\star)$, target unemployment rate $u^\star$}
\KwOut{Calibrated $c^\star$, equilibrium objects $(\theta^\star,w^\star,q^\star,f^\star,u^\star)$}
\BlankLine
\Repeat{$|u - u^\star| < \varepsilon_\kappa$}{
  \tcp{Inner steady-state problem given $c$}
  Use the bisection (or alternative) algorithm\;
  \Repeat{$|q^{n+1} - q^n| < \varepsilon_q$}{
    Compute implied tightness $\theta$ from Equation \eqref{jfr} \;
    Compute implied wage $w$ from Equation \eqref{eq_wage} \;
    Compute $\tilde\theta$ from the job-creation condition, Equation \eqref{jobcreation}  \;
    Update $q$ by bisection on the residual $|q^{n+1} - q^n|$\;
  }
  \tcp{Evaluate calibration target}
  Compute job-finding rate $f$ from the matching function\;
  Compute steady-state unemployment $u$ from Equation \eqref{beveridge}\;
  Update $c$ by bisection on the residual $u - u^\star$\;
}
\end{algorithm}

\printbibliography
\end{document}


%------------------------------------------------------------------------
%Useful Snippets.
%------------------------------------------------------------------------
%\section{Own Snippets (Comment Out)}
%\vspace{3mm}
%
%% 1. Code Listing. 
%\begin{lstlisting}[frame=single,language=Matlab]
%	An example of Matlab code listing.      
%\end{lstlisting}
%
%% 2. Highlighting.
%\colorbox{maroon!10}{An example of highlighted text.}	
%
%%3. Two Figures Next to Each Other.
%\begin{figure}[hbt]
%\centering
%\begin{subfigure}{.5\textwidth}
%  \centering
%  \includegraphics[width=\linewidth]{}
%  \caption{A}
%\end{subfigure}%
%\begin{subfigure}{.5\textwidth}
%  \centering
%  \includegraphics[width=\linewidth]{}
%  \caption{B}
%\end{subfigure}
%\caption{Two Figures Next to Each Other}
%\label{two_figures}
%\end{figure}	

